\chapter{Discusión de los Resultados}
En este capítulo se someten a un análisis crítico los resultados empíricos expuestos en el capítulo previo. El propósito fundamental es evaluar la eficiencia, viabilidad y éxito del diseño e implementación de las técnicas de poda a nivel de nodo hoja sobre la arquitectura del Octree lineal en relación a los resultados y soluciones existentes a día de hoy para la búsqueda de vecindades en nubes de puntos capturadas con tecnología LiDAR.

\section{Análisis de Estructuras de Almacenamiento: algoritmo \texttt{neighborsStruct} frente a \texttt{neighborsPrune}}

Como tendencia independiente de las contribuciones directas de este trabajo, se constata que la arquitectura de almacenamiento de puntos vecinos basado en pares de índices, (\textit{neighborsStruct}), reporta de manera sistemática una eficiencia temporal superior a las inserciones secuenciales (presentes en \textit{neighborsPrune}). Este comportamiento se explica por la reducción de inserciones en el vector \textit{ptsInside} realizadas, diferencia que se acentúa más cuantos más puntos consecutivos sean considerados válidos.

Sin embargo, el hallazgo más relevante de esta comparativa radica en el comportamiento asimétrico que manifiesta el algoritmo de poda polar. Esta optimización exhibe un impacto sustancialmente mayor al acoplarse sobre \textit{neighborsStruct} que al operar sobre \textit{neighborsPrune}. La hipótesis que lo justifica reside en el modo de inserción en el vector de resultados:

\begin{itemize}
    \item En la variante básica \textit{neighborsPrune}, el algoritmo evalúa secuencialmente los puntos y añade un elemento al contenedor de salida de forma individualizada si, y solo si, se verifica la condición de vecindad positiva. Bajo este esquema, la poda angular actúa únicamente reduciendo el coste de la evaluación \textit{booleana} interna (\texttt{k.isInside()}), dado que todos los puntos descartados habrían devuelto un valor falso en cualquier caso. El ahorro computacional queda limitado a la supresión de la comprobación aritmética de \texttt{isInside()}.
    
    \item Por el contrario, en \textit{neighborsStruct} la lógica opera por bloques de rangos contiguos mediante la guarda condicional \texttt{if (!k.isInside(points[i]))}. Si un punto no pertenece al entorno, se introduce en el contenedor el rango completo indexado previamente. Al efectuar la poda a nivel de hoja, cuanto más efectiva sea la discriminación geométrica, menor será el número de ocasiones en las que el flujo de ejecución entra en esta bifurcación, reduciendo masivamente las operaciones de inserción y de redimensionado del vector de salida.
\end{itemize}

Esta divergencia es particularmente visible bajo kernels esféricos, ya que es la forma geométrica adaptada al cálculo de poda angular que se realiza en el modo polar. Aparte, en nubes de puntos con una cierta densidad espacial (tales como \textit{Lille\_0}, \textit{bildstein\_station1}, \textit{sg27\_station8} o \textit{Mar18\_test}), se aprecia que a medida que se incrementa el radio operativo, el número de puntos consecutivos que cumplen con el criterio de inclusión se dispara. En estas regiones de alta saturación, el algoritmo de pares de índices de \textit{neighborsStruct} amplía la diferencia proporcionalmente en rendimiento contra el \textit{pruning} estándar.

\section{Sensibilidad al Tamaño de Hoja y Dinámicas del Kernel}

Un hito de la arquitectura propuesta es su capacidad para romper la degradación del rendimiento asociada al crecimiento del tamaño de los nodos hoja (\texttt{maxLeaf}). En las implementaciones secuenciales clásicas, incrementar la capacidad del nodo hoja implica obligatoriamente dilatar el tiempo de cómputo, al verse forzado el sistema a realizar un escaneo lineal sobre un mayor volumen de puntos desordenados.

Con la introducción de los modos de reordenamiento local, el coste temporal deja de depender directamente del número absoluto de puntos contenidos en la hoja. El rendimiento pasa a depender de la distribución espacial y del porcentaje de descarte geométrico efectivo. Este porcentaje de descarte fue estudiado con detalle en el capítulo anterior, y se pueden extraer una serie de aprendizajes sobre el funcionamiento teórico, más allá de los tiempos de ejecución producidos. El reordenamiento local polar funciona mejor en combinación con el modo de  kernel esférico por la concordancia geométrica existente, mientras que para el caso de los cubos hay que trazar una esfera que lo circunscriba para poder hallar esos límites. Además, también se entiende la mayor efectividad de poda en radios bajos. La sección angular mínima que contiene a la esfera de búsqueda es más cerrada cuánto más pequeña sea esa esfera. Esta tendencia esperada sí se corresponde con los resultados del algoritmo \textit{neighborsStruct} en las Figuras \ref{fig:subset_lines_train} y \ref{fig:subset_lines_test}. De forma complementaria, la efectividad de la poda en los experimentos estudiados para el reordenamiento cartesiano tuvo éxito en las búsquedas de geometría cúbica, al separar perfectamente los puntos según los ejes ordenados. En estos casos, se hallaba una dinámica contrapuesta al variar el tamaño del radio: las podas son más débiles en radios pequeños. Esto tiene una explicación geométrica detrás, y es que en búsquedas cúbicas, si el kernel solo corta la hoja en una dimensión la separación de puntos entre vecinos y no vecinos será perfecta. Solo en los casos en los que dos o tres dimensiones cortan el nodo se producirán podas imperfectas. En cubos grandes (proporcionalmente a los tamaños de los nodos terminales) más octantes tendrán la cualidad de ser cortados solamente por una dimensión, con lo que más puntos serán separados de forma perfecta. Sin embargo, este comportamiento no se aprecia directamente en los experimentos realizados usando el algoritmo \textit{neighborsPrune}, y tiene una explicación que se introducirá a continuación.


\section{Evaluación de los Modos de Reordenamiento Local}

Al contrastar de manera aislada los dos modos de reordenamiento implementados (polar y cartesiano) frente al algoritmo de búsqueda con inserciones secuenciales, las métricas reflejan diferencias muy sutiles. Esto se debe a que existe un beneficio extra en que se inserten una serie de puntos consecutivos, como sí pasa al usar \textit{neighborsStruct}. A pesar de ello, la introducción de la poda logra amortizar con éxito su propia sobrecarga algorítmica, manteniendo registros competitivos.

Teóricamente, cabría esperar una divergencia de rendimiento más acusada entre ambas reordenaciones al cruzar sus kernels respectivos, dado que cada una minimiza la sobreestimación del volumen envolvente en su propio dominio geométrico. Sin embargo, la mejora temporal respecto a la versión secuencial es tan pequeña que el ahorro funcional apenas sobresale respecto al coste fijo de computar los límites de la hoja. Hay casos concretos en los que, tanto el modo polar como el cartesiano, ganan rendimiento al ser acoplados sobre \textit{neighborsPrune}, pero no de forma consistente ni con una tendencia clara.

En este contexto en el que la sobrecarga prácticamente alcanza al ahorro conseguido mediante la poda, los escenarios cruzados (modo cartesiano sobre kernels esféricos y modo polar sobre kernels cúbicos) operan como un estado intermedio que tampoco degrada apenas el rendimiento. En estos casos, la sobreestimación geométrica del volumen impide que la mayoría de las hojas califiquen para un descarte. No obstante, estas hojas no aptas evitan la ejecución de la búsqueda binaria y se desvían de forma inmediata hacia el recorrido secuencial ordinario. El software no penaliza de forma crítica estas configuraciones subóptimas de entrada, garantizando la estabilidad del rendimiento.

\section{Calibración del Espacio de Hiperparámetros}

Los ensayos demuestran que el rendimiento de las consultas espaciales es altamente sensible al ecosistema de parámetros de control, condicionado de manera simultánea por la densidad de la nube, la magnitud del radio y la posición espacial del centro de consulta. 

A pesar de esta volatilidad, es posible extraer una regla de diseño generalizada: la configuración de tamaños de nodo hoja superiores a \texttt{maxLeaf = 256} introduce un principio de degradación en la eficiencia temporal. Esta pérdida de rendimiento, cabe destacar, es significativamente atenuada en comparación con la bajada de rendimiento en los algoritmos clásicos. Aún así, delimita el umbral que no se debería superar a menos que se esté trabajando sobre un entorno ya conocido y estudiado.

En lo que respecta al umbral mínimo de puntos para la activación de la poda (\texttt{threshold}), el valor de \texttt{64} puntos se consolida como el punto de equilibrio óptimo. Límites inferiores introducen un \textit{overhead} ineficiente en la CPU al forzar el cálculo de selectores de rango para subconjuntos de puntos minúsculos cuya evaluación secuencial directa en memoria sería más veloz. Por el contrario, límites excesivamente elevados clausuran ventanas de oportunidad de la poda en hojas medianas que albergan buen potencial de descarte.

\section{Dinámicas en Búsquedas de Cobertura Completa}

El escenario de evaluación sobre búsquedas completas y secuenciales (\textit{Full Searches}) proporcionó un escenario de resultados muy dispar a los discutidos hasta ahora. En este modo de exploración, la tasa de fallos de caché se reduce al mínimo debido a la asignación por bloques de puntos a cada hilo OpenMP.

Al procesarse núcleos de consulta vecinos de forma consecutiva, hay una gran tasa de reutilización de puntos en los niveles superiores de la memoria caché de la CPU tras las iteraciones previas. Este fenómeno se amplifica teniendo en cuenta que los puntos fueron reordenados de forma global mediante las curvas de Hilbert para mejorar su localidad. Como consecuencia directa de esta localidad hardware, el coste temporal de evaluar secuencialmente un conjunto de puntos que potencialmente se iban a descartar es extremadamente bajo. 

Por ende, dado que el coste de cálculo del \textit{overhead} para la selección de rangos geométricos se mantiene constante, el balance de rentabilidad de la poda es negativo. Los tiempos de ejecución totales, aunque no suponen una degradación notable, se sitúan ligeramente por detrás de las búsquedas con los algoritmos originales. En este contexto de coberturas completas, no se considera útil la implementación actual de reordenamientos locales, ya que no produce un rendimiento consistentemente mejor al código de partida a la vez que otorga a los \textit{Octrees} un recorrido de búsqueda más complejo y produce un mayor gasto en memoria. 